\documentclass[leqno]{article}
% Shared mark scheme preamble for resources.danieljsmith.org: the pack
% preamble, unchanged, then the scheme's own layer. A scheme's .tex is
% rendered from its YAML (scripts/README.md, Mark schemes), never edited.
\input{\CurrentFilePath/pack_preamble.tex}

% The page: equal side margins, so the table runs the full width of the
% header rule. The footers' right edge follows the text's.
\geometry{left=0.8in, right=0.8in}
\makeatletter
\renewcommand{\djs@footR}[1]{%
  \put(\dimexpr1in+\hoffset+\oddsidemargin+\textwidth\relax,
    -\dimexpr\paperheight-0.32in\relax){\makebox[0pt][r]{#1}}}
\makeatother

\usepackage[export]{adjustbox}
\usepackage{tabularray}
\definecolor{msGrid}{RGB}{176, 172, 164}
\definecolor{msTint}{RGB}{245, 243, 239}

% The header on every page, the front page included, and the booklet label
% under the contents-page title. \djsQuestionPageStyle stays empty.
\newcommand{\djsSchemeHeader}[1]{%
  \djsPackHeader[markscheme]{\textit{Mark scheme} - #1}%
  \renewcommand{\djsPackTitle}{#1}%
  \renewcommand{\djsBookletLabel}{Mark scheme}%
}

% The front page: the contents heading, spaced as tocloft spaces the packs'
% "Questions", then the contents table. A question's name and page number
% both link to its first page.
\newcommand{\msContentsHeading}{%
  \vskip\cftbeforetoctitleskip
  {\noindent\cfttoctitlefont\contentsname\par}%
  \vskip\cftaftertoctitleskip}
\NewTblrEnviron{mscontents}
\SetTblrInner[mscontents]{
  colspec = {Q[wd=26mm, l] X[l] Q[wd=16mm, c] Q[wd=14mm, r]},
  width = \linewidth,
  rowsep = 4pt, colsep = 6pt,
  hlines = {0.4pt, msGrid}, vlines = {0.4pt, msGrid},
  hline{2} = {2.2pt, scarlet},
  row{1} = {font=\small\bfseries},
  row{Z} = {font=\bfseries},
  column{2} = {font=\small},
}
\newcommand{\msContentsQuestion}[1]{\hyperref[q:#1]{Question #1}}
\newcommand{\msContentsPage}[1]{\hyperref[q:#1]{\pageref*{q:#1}}}
\newcommand{\msSep}{\hspace{0.9em}}

% The conventions pages' layout: headings in the documents' own bold, tables
% ruled and padded as the scheme's, roomy lists.
\newcommand{\msHead}[1]{%
  \par\vspace{18pt}%
  {\noindent\bfseries #1\par}%
  \vspace{6pt}}
\NewTblrEnviron{mstblr}
\SetTblrInner[mstblr]{rowsep = 6pt, colsep = 6pt,
  hlines = {0.4pt, msGrid}, vlines = {0.4pt, msGrid}}
\newenvironment{msList}
  {\begin{itemize}[leftmargin=1.2em, itemsep=5pt, topsep=3pt, parsep=0pt]}
  {\end{itemize}}

% A question's first page: the link target and the PDF bookmark, as
% \questionitem gives the packs.
\newcommand{\msQnum}{}
\newcommand{\msQuestion}[1]{%
  \clearpage
  \renewcommand{\msQnum}{#1}%
  \phantomsection\label{q:#1}%
  \addcontentsline{toc}{subsection}{\protect\djsTocEntry{#1}}}

% The scheme's table: Question 16mm, Scheme the rest, Marks 16mm, Notes
% 44mm; every cell ruled in msGrid; the scarlet rule under the head row. A
% long table has no caption or "continued" foot; a continuation page opens
% with a muted line above the repeated head row. The rules across a part's
% rows, the alternative's tint and the total row are the rendered table's.
\NewTblrTheme{ms}{%
  \DefTblrTemplate{firsthead}{default}{}%
  \DefTblrTemplate{middlehead,lasthead}{default}{%
    {\small\itshape\color{djsMuted}Question \msQnum\ continued\par}%
    \vspace{6pt}}%
  \DefTblrTemplate{firstfoot,middlefoot,lastfoot}{default}{}%
}
\NewTblrEnviron{msscheme}
\SetTblrOuter[msscheme]{long, theme=ms, entry=none, label=none}
\SetTblrInner[msscheme]{
  colspec = {Q[wd=16mm, l] X[l] Q[wd=16mm, c] Q[wd=44mm, l]},
  width = \linewidth,
  rowsep = 6pt, colsep = 6pt,
  rowhead = 1,
  vlines = {0.4pt, msGrid},
  hline{1} = {0.4pt, msGrid}, hline{2} = {2.2pt, scarlet},
  hline{Z} = {0.4pt, msGrid},
  column{1} = {font=\bfseries}, column{3} = {font=\bfseries},
  column{4} = {font=\small},
  row{1} = {rowsep=4pt, font=\small\bfseries},
}

% Two marks on a row: each note starts with its code.
\newcommand{\msNote}[2]{\textbf{#1}\enspace #2}
% A full-row figure: the worked solution's own picture, scaled uniformly and
% only down, to the table less the Question column. The picture is drawn
% with the line width it has in the solutions document, at its list depth
% there (the optional argument, 1 by default: a part's working), so one
% sized by \linewidth keeps its size; tabularray sets the cell's \linewidth
% to the cell's measured width, which would shrink it at every pass. The
% widths are measured once, in a solution box at the packs' text width
% (pack_preamble.tex's margins, 0.8in and 1in).
\newlength{\msSpanW}
\makeatletter
\newcommand{\ms@keepLW}[1]{%
  \expandafter\xdef\csname ms@solnLW@#1\endcsname{\the\linewidth}}
\AtBeginDocument{%
  \setlength{\msSpanW}{\dimexpr\textwidth-16mm-28pt\relax}%
  \setbox\z@\vbox{%
    \begin{minipage}{\dimexpr\paperwidth-1.8in\relax}
      \begin{djsSolution}
        \ms@keepLW{0}%
        \begin{enumerate}\item \ms@keepLW{1}%
        \begin{enumerate}\item \ms@keepLW{2}%
        \begin{enumerate}\item \ms@keepLW{3}%
        \end{enumerate}\end{enumerate}\end{enumerate}
      \end{djsSolution}
    \end{minipage}}}
\newcommand{\msFigure}[2][1]{%
  \adjustbox{max width=\msSpanW, max totalheight=0.6\textheight, valign=t}{%
    \hsize=\csname ms@solnLW@#1\endcsname
    \linewidth=\hsize
    #2}}
\makeatother
% A figure mark's feature, in the Scheme column.
\newcommand{\msFeature}[1]{\textit{#1}}
% An alternative's title row.
\newcommand{\msAlternative}[1]{\textit{Alternative:} #1}
% Drafts only: a queried mark's code, the query's why under its note, and a
% flag in its part's place.
\newcommand{\msQueried}[2]{#1~\(\to\) #2?}
\newcommand{\msQuery}[1]{{\itshape\color{djsMuted}Query: #1}}
\newcommand{\msFlag}[2]{\textbf{#1:}\enspace #2}


\djsSchemeHeader{Proof by Induction - Matrices}
\djsSitePath{Core Pure}{Proof by Induction}

\djsLastUpdated{25 October 2026}
\begin{document}
\djsSetQuestionCount{8}
\djsAddMarks{1}{6}
\djsAddMarks{2}{6}
\djsAddMarks{3}{8}
\djsAddMarks{4}{6}
\djsAddMarks{5}{11}
\djsAddMarks{6}{6}
\djsAddMarks{7}{6}
\djsAddMarks{8}{6}
\vspace*{6pt}
\djsContentsTitle
\vspace{30pt}
\msContentsHeading
\begin{mscontents}{}
Question & Marks by part & Marks & Page \\
\msContentsQuestion{1} & & 6 & \msContentsPage{1} \\
\msContentsQuestion{2} & & 6 & \msContentsPage{2} \\
\msContentsQuestion{3} & (a)\,6\msSep (b)\,2 & 8 & \msContentsPage{3} \\
\msContentsQuestion{4} & & 6 & \msContentsPage{4} \\
\msContentsQuestion{5} & (a)\,6\msSep (b)\,5 & 11 & \msContentsPage{5} \\
\msContentsQuestion{6} & & 6 & \msContentsPage{6} \\
\msContentsQuestion{7} & & 6 & \msContentsPage{7} \\
\msContentsQuestion{8} & & 6 & \msContentsPage{8} \\
\SetCell[c=2]{r} Total & & 55 & \\
\end{mscontents}
\newpage
% The conventions pages, shared by every pack's mark scheme: one source of
% truth, \input after the front page, with no strand or pack logic. The keys
% on the first side, the rules on the second. The scheme layer defines
% \msHead, msList and mstblr.
\vspace*{6pt}%
{\noindent\cfttoctitlefont Reading the mark scheme\par}%
\vspace{10pt}%
These rules apply throughout, unless a note says otherwise.

\msHead{Marks}
Every mark is worth one, so \textbf{M1 A1} on a row is two marks. Only marks
coded ft follow through.\par
\vspace{10pt}%
\begin{mstblr}{
  colspec = {Q[wd=15mm, l] X[l]},
  width = \linewidth,
  hline{2} = {2.2pt, scarlet},
  column{1} = {font=\bfseries},
  row{1} = {font=\small\bfseries, rowsep=4pt},
}
Code & Meaning \\
M1 & Method: a correct method, applied to this question. A slip doesn't cost it if the method is clear. A formula must be used, not just quoted. \\
dM1 & Dependent method: as M1, but only if the method mark before it was earned, unless its note names another. \\
A1 & Accuracy: a correct answer or step. It needs the method mark it follows. \\
A1ft & As A1, following through from an earlier answer that the note names. \\
A1~AG & The last mark of a part whose answer is printed in the question: complete working, with no errors. \\
B1 & A correct value, fact or statement. It needs no method mark. \\
B1ft & As B1, following through from an earlier answer that the note names. \\
E1 & A written explanation, interpretation or conclusion. Concluding a proof by induction needs every earlier mark. \\
\end{mstblr}

\msHead{Abbreviations}
ft and AG qualify a code in the Marks column. The others appear in the
notes.\par
\vspace{10pt}%
\begin{mstblr}{
  colspec = {Q[wd=17mm, l] X[l] Q[wd=15mm, l] X[l]},
  width = \linewidth,
  column{1} = {font=\bfseries}, column{3} = {font=\bfseries},
}
ft & follow through & AG & answer given in the question \\
cao & correct answer only, with no follow-through & cso & correct solution only: no errors in the part \\
oe & or an equivalent form & awrt & any answer that rounds to this value \\
soi & seen, or implied by later working & isw & ignore working after a correct answer \\
SC & special case: credit outside the main method & their & the student's own earlier value \\
condone & a slip that doesn't cost the mark & [\;\,] & working the mark doesn't need \\
s.f. & significant figures & d.p. & decimal places \\
\end{mstblr}

\newpage
\msHead{Reading the scheme}
\begin{msList}
\item Each row is a line of working, with the marks it earns beside it.
  Where a row carries two marks, each note starts with its code.
\item \textbf{[\textit{n}]} ends each part with its total.
\item A shaded \textit{Alternative} is a second method for the same total.
  Mark an answer by the method it follows.
\item Any other valid method earns the equivalent marks, unless the question
  names a method or says \textit{hence}.
\item A sketch earns one B1 for each feature named below its figure.
\end{msList}

\msHead{Accuracy and answers}
\begin{msList}
\item Unless a question says otherwise, give answers to 3 significant
  figures, and angles in degrees to 1 decimal place. Angles in radians take
  3 significant figures.
\item A final answer with more figures is accepted if it rounds to the
  scheme's value; one with fewer is not. An answer spoiled by rounding too
  early loses its accuracy mark.
\item Where a question asks for an exact answer, a decimal loses the final
  accuracy mark. Otherwise either is accepted.
\item Wrong or missing units cost nothing unless the question asks for them.
\item A correct answer with no working earns full marks, unless the question
  needs the working: show that, prove, hence, a method it names, or an exact
  answer. A wrong answer with no working earns nothing. Calculator answers
  are treated the same way.
\item A correct answer, once reached, isn't spoiled by later working (isw).
\end{msList}

\msHead{Value of $\boldsymbol{g}$}
\begin{msList}
\item Take $g = 9.8\,\mathrm{m\,s^{-2}}$ unless a question states a value,
  which is then the only one accepted. Answers that use $g = 9.8$ may be
  given to 2 or 3 significant figures.
\item Where a question doesn't state $g$, an answer worked with $g = 10$
  throughout earns the same marks, keeping any value printed in the
  question. The worked solutions list the final answers with $g = 10$.
\end{msList}


\msQuestion{1}
\begin{msscheme}{
  hline{3-6} = {2-Z}{0.4pt, msGrid},
  hline{7} = {0.4pt, msGrid},
  row{7} = {font=\bfseries},
}
Question & Scheme & Marks & Notes \\
1 & {\(n=1\): \(\begin{pmatrix}2^1 & 3^1-2^1\\0 & 3^1\end{pmatrix} = \begin{pmatrix}2 & 1\\0 & 3\end{pmatrix}\)} & B1 & {Needs the formula evaluated at \(n=1\) and shown equal to the given matrix; 'true for \(n=1\)' alone is not enough.} \\
& {Assume true for \(n=k\):\\[4pt] \(\begin{pmatrix}2 & 1\\0 & 3\end{pmatrix}^{k+1}\)\\[4pt] \(\phantom{A} = \begin{pmatrix}2^k & 3^k-2^k\\0 & 3^k\end{pmatrix}\begin{pmatrix}2 & 1\\0 & 3\end{pmatrix}\)} & M1 & {States the assumption for \(n=k\) and writes the \((k+1)\)th power as the assumed \(k\)th-power matrix times the given matrix, in either order.} \\
& {\(\phantom{A} = \begin{pmatrix}2^{k+1} & 2^k+3(3^k-2^k)\\0 & 3^{k+1}\end{pmatrix}\)} & M1 A1 & {\msNote{M1}{Multiplies their matrices in either order, forming the top-right entry row by column. Condone slips.}\\ \msNote{A1}{Correct product in either order; entries need not be simplified.}} \\
& {\(2^k+3^{k+1}-3\times2^k = 3^{k+1}-2^{k+1}\)} & A1 & {Top-right entry correctly simplified to \(3^{k+1}-2^{k+1}\), other entries \(2^{k+1}\), \(0\), \(3^{k+1}\). Not for the result asserted directly from the product.} \\
& {True for \(n=1\); true for \(n=k\)\\[4pt] \(\Rightarrow\) true for \(n=k+1\); so true for all \(n\in\mathbb{Z}^+\)} & E1 & {Concludes with all three points: true for \(n=1\), true for \(n=k\) implies true for \(n=k+1\), hence true for all positive integers.} \\
\nopagebreak \SetCell[c=2]{r} Total for Question 1 & & 6 & \\
\end{msscheme}

\msQuestion{2}
\begin{msscheme}{
  hline{3-6} = {2-Z}{0.4pt, msGrid},
  hline{7} = {0.4pt, msGrid},
  row{7} = {font=\bfseries},
}
Question & Scheme & Marks & Notes \\
2 & {\(n=1\): LHS \(= \begin{pmatrix}3 & -1\\ 1 & 1\end{pmatrix}\)\\[4pt] RHS \(= 2^{0}\begin{pmatrix}3 & -1\\ 1 & 1\end{pmatrix} = \begin{pmatrix}3 & -1\\ 1 & 1\end{pmatrix}\)} & B1 & {Both sides evaluated at \(n=1\) and shown equal; a bare statement that it is true scores B0.} \\
& {Assumes \(\mathbf{A}^k = 2^{k-1}\begin{pmatrix}k+2 & -k\\ k & 2-k\end{pmatrix}\)\\[4pt] \(\mathbf{A}^{k+1} = 2^{k-1}\begin{pmatrix}k+2 & -k\\ k & 2-k\end{pmatrix}\begin{pmatrix}3 & -1\\ 1 & 1\end{pmatrix}\)} & M1 & {Assumes the result for \(n=k\) and writes \(\mathbf{A}^{k+1}\) as their \(\mathbf{A}^k\) times \(\mathbf{A}\), in either order.} \\
& {\(\phantom{\mathbf{A}^{k+1}} = 2^{k-1}\begin{pmatrix}3(k+2)-k & -(k+2)-k\\ 3k+2-k & -k+2-k\end{pmatrix}\)} & M1 A1 & {\msNote{M1}{Multiplies the matrices in the order set up, attempting all four elements. Condone one slip.}\\ \msNote{A1}{Correct product for their order, need not be simplified.}} \\
& {\(\phantom{\mathbf{A}^{k+1}} = 2^{k}\begin{pmatrix}k+3 & -(k+1)\\ k+1 & 2-(k+1)\end{pmatrix}\)} & A1 & {cso. Takes out the factor 2 to reach the result with \(n=k+1\), shown in that form or stated as the \(n=k+1\) result.} \\
& {True for \(n=1\); true for \(n=k\) \(\Rightarrow\) true for \(n=k+1\)\\[4pt] \(\Rightarrow\) true for all \(n \in \mathbb{N}\)} & E1 & {States true for \(n=1\), true for \(n=k\) implies true for \(n=k+1\), and hence true for all \(n \in \mathbb{N}\).} \\
\nopagebreak \SetCell[c=2]{r} Total for Question 2 & & 6 & \\
\end{msscheme}

\msQuestion{3}
\begin{msscheme}{
  hline{3-7,9-10} = {2-Z}{0.4pt, msGrid},
  hline{8,11} = {0.4pt, msGrid},
  row{11} = {font=\bfseries},
}
Question & Scheme & Marks & Notes \\
3(a) & {\(n=1\): \(2^1\begin{pmatrix}2&1\\-1&0\end{pmatrix}=\begin{pmatrix}4&2\\-2&0\end{pmatrix}=\mathbf{A}\)} & B1 & {Formula evaluated at \(n=1\) and shown equal to \(\mathbf{A}\); a bare statement that it is true for \(n=1\) scores nothing.} \\
& {Assume \(\mathbf{A}^k=2^k\begin{pmatrix}k+1&k\\-k&1-k\end{pmatrix}\)\\[4pt] \(\mathbf{A}^{k+1}=2^k\begin{pmatrix}k+1&k\\-k&1-k\end{pmatrix}\begin{pmatrix}4&2\\-2&0\end{pmatrix}\)} & M1 & {States the assumption for \(n=k\) and writes \(\mathbf{A}^{k+1}\) as \(\mathbf{A}^k\mathbf{A}\) or \(\mathbf{A}\mathbf{A}^k\) with the assumed form substituted.} \\
& {\(\phantom{\mathbf{A}^{k+1}}=2^k\begin{pmatrix}4(k+1)-2k&2(k+1)\\-4k-2(1-k)&-2k\end{pmatrix}\)} & M1 A1 & {\msNote{M1}{Multiplies row by column for all four elements, in either order. Condone arithmetic slips.}\\ \msNote{A1}{Correct product in either order; need not be simplified.}} \\
& {\(\phantom{\mathbf{A}^{k+1}}=2^{k+1}\begin{pmatrix}k+2&k+1\\-(k+1)&1-(k+1)\end{pmatrix}\)} & A1 & {Factor \(2^{k+1}\) with every element shown equal to the formula's with \(n=k+1\), by rewriting or by comparison with a stated target.} \\
& {True for \(n=1\); true for \(n=k\) \(\Rightarrow\) true for \(n=k+1\)\\[4pt] \(\Rightarrow\) true for all positive integers \(n\)} & E1 & {States true for \(n=1\), true for \(n=k\) implies true for \(n=k+1\), hence true for all positive integers \(n\).} \\
\nopagebreak & & [6] & \\
3(b) & {\(\mathbf{A}^n=2^n\mathbf{M}\), \(\mathbf{M}=\begin{pmatrix}n+1&n\\-n&1-n\end{pmatrix}\)\\[4pt] \(\det\mathbf{M}=(n+1)(1-n)+n^2=1\)\\[4pt] \(\mathbf{B}=2^{-n}\mathbf{M}^{-1}\)} & M1 & {Takes out \(2^n\), or works with \(\mathbf{A}^n\): swaps the leading diagonal, negates the other two elements and divides by their determinant. Condone determinant slips.} \\
& {\(\mathbf{B}=\frac{1}{2^n}\begin{pmatrix}1-n&-n\\n&n+1\end{pmatrix}\)} & A1 & {oe, such as \(2^{-n}\) or the factor multiplied into each element.} \\
\nopagebreak & & [2] & \\
\nopagebreak \SetCell[c=2]{r} Total for Question 3 & & 8 & \\
\end{msscheme}

\msQuestion{4}
\begin{msscheme}{
  hline{3-6} = {2-Z}{0.4pt, msGrid},
  hline{7} = {0.4pt, msGrid},
  row{7} = {font=\bfseries},
}
Question & Scheme & Marks & Notes \\
4 & {\(n = 1\): \(2^1 = 2\), \(1\times 2^0 = 1\), \(\frac{1(0)}{2}2^{-1} = 0\)\\[4pt] \(\Rightarrow\) RHS \(= \mathbf{A}\)} & B1 & {RHS at \(n = 1\) evaluated and shown equal to \(\mathbf{A}\); not for stating "true for \(n=1\)" alone.} \\
& {Assume true for \(n = k\): \(\mathbf{A}^k = \) the given matrix with \(n = k\)\\[4pt] \(\mathbf{A}^{k+1} = \mathbf{A}^k\mathbf{A}\)} & M1 & {States the assumption for \(n = k\) and writes \(\mathbf{A}^{k+1}\) as \(\mathbf{A}^k\mathbf{A}\) or \(\mathbf{A}\mathbf{A}^k\) with the assumed matrix.} \\
& {Diagonal: \(2^k\times 2 = 2^{k+1}\)\\[4pt] Below diagonal: \(k2^{k-1}\times 2 + 2^k\times 1\)\\[4pt] Bottom left: \(\frac{k(k-1)}{2}2^{k-2}\times 2 + k2^{k-1}\times 1\)} & M1 A1 & {\msNote{M1}{Multiplies the two matrices, attempting at least the entries below the diagonal. Condone slips.}\\ \msNote{A1}{All nine entries correct, unsimplified.}} \\
& {\(k2^k + 2^k = (k+1)2^k\)\\[4pt] \(\left(\frac{k(k-1)}{2}+k\right)2^{k-1} = \frac{k(k+1)}{2}2^{k-1}\)\\[4pt] \(\phantom{\left(\frac{k(k-1)}{2}+k\right)2^{k-1}} = \frac{(k+1)k}{2}2^{(k+1)-2}\)} & A1 & {cso. Every entry in the \(n=k+1\) form, the bottom-left reaching \(\frac{k(k+1)}{2}2^{k-1}\) oe by correct algebra.} \\
& {True for \(n=1\); true for \(n=k\) \(\Rightarrow\) true for \(n=k+1\)\\[4pt] \(\Rightarrow\) true for all \(n \geq 1\)} & E1 & {Concludes with all three: true for \(n=1\); true for \(n=k\) implies true for \(n=k+1\); hence true for all integers \(n \geq 1\).} \\
\nopagebreak \SetCell[c=2]{r} Total for Question 4 & & 6 & \\
\end{msscheme}

\msQuestion{5}
\begin{msscheme}{
  hline{3-7,9-12} = {2-Z}{0.4pt, msGrid},
  hline{8,13} = {0.4pt, msGrid},
  row{13} = {font=\bfseries},
}
Question & Scheme & Marks & Notes \\
5(a) & {\(n=1\): \(\begin{pmatrix}(3c)^1 & 1(3c)^0\\0&(3c)^1\end{pmatrix} = \begin{pmatrix}3c&1\\0&3c\end{pmatrix} = \mathbf{M}\)} & B1 & {Needs the RHS with \(n=1\) substituted shown equal to \(\mathbf{M}\).} \\
& {Assume \(\mathbf{M}^k=\begin{pmatrix}(3c)^k & k(3c)^{k-1}\\0&(3c)^k\end{pmatrix}\)\\[4pt] \(\mathbf{M}^{k+1}=\mathbf{M}^k\mathbf{M}\)} & M1 & {States the assumption for \(n=k\) and writes \(\mathbf{M}^{k+1}\) as \(\mathbf{M}^k\mathbf{M}\) or \(\mathbf{M}\mathbf{M}^k\).} \\
& {\(\begin{pmatrix}(3c)^k & k(3c)^{k-1}\\0&(3c)^k\end{pmatrix}\begin{pmatrix}3c&1\\0&3c\end{pmatrix}\)\\[4pt] \(\phantom{A}=\begin{pmatrix}(3c)^{k+1} & (3c)^k+3ck(3c)^{k-1}\\0&(3c)^{k+1}\end{pmatrix}\)} & M1 A1 & {\msNote{M1}{Multiplies the assumed \(\mathbf{M}^k\) and \(\mathbf{M}\), in either order, by rows and columns. Condone slips in entries.}\\ \msNote{A1}{Correct product, any unsimplified form.}} \\
& {\(\phantom{A}=\begin{pmatrix}(3c)^{k+1} & (k+1)(3c)^k\\0&(3c)^{k+1}\end{pmatrix}\)} & A1 & {Top-right entry reaches \((k+1)(3c)^k\), oe such as \((k+1)(3c)^{(k+1)-1}\), without error.} \\
& {True for \(n=1\); true for \(n=k\) \(\Rightarrow\) true for \(n=k+1\);\\[4pt] so true for all positive integers \(n\)} & E1 & {States the base case, the implication from \(n=k\) to \(n=k+1\), and truth for all positive integers \(n\).} \\
\nopagebreak & & [6] & \\
5(b) & {\(\det(\mathbf{M}^n)=(3c)^n(3c)^n-n(3c)^{n-1}\times 0\)} & M1 & {Uses \(ad-bc\) on the matrix from (a), or \((\det\mathbf{M})^n\).} \\
& {\(\phantom{\det(\mathbf{M}^n)}=(3c)^{2n}\)} & A1 & {oe, such as \((9c^2)^n\) or \(9^nc^{2n}\).} \\
& {\((9c^2)^n=144^n \Rightarrow 9c^2=144\)} & M1 A1 & {\msNote{M1}{Equates their determinant, of the form \((kc^2)^n\) or \((kc)^{2n}\), to \(144^n\) and removes the power \(n\) by roots, bases or logs.}\\ \msNote{A1}{oe, such as \(3c=12\) or \(c^2=16\).}} \\
& {\(c=4\)} & A1 & {A0 if \(c=-4\) is also given as an answer.} \\
\nopagebreak & & [5] & \\
\nopagebreak \SetCell[c=2]{r} Total for Question 5 & & 11 & \\
\end{msscheme}

\msQuestion{6}
\begin{msscheme}{
  hline{3-6} = {2-Z}{0.4pt, msGrid},
  hline{7} = {0.4pt, msGrid},
  row{7} = {font=\bfseries},
}
Question & Scheme & Marks & Notes \\
6 & {\(n=1\): \((-1)^1 = -1\), \(-1(-1)^1 = 1\), \(\frac{1(0)}{2}(-1)^1 = 0\)\\[4pt] RHS = given matrix, so true for \(n=1\)} & B1 & {Evaluates the RHS at \(n=1\) and states it equals the given matrix. Writing the given matrix alone is not enough.} \\
& {Assumes true for \(n=k\); writes the \((k+1)\)th power\\[4pt] as the \(k\)th-power matrix times the given matrix} & M1 & {States the assumption for \(n=k\) and forms the product of the assumed matrix and the given matrix, in either order.} \\
& {Leading diagonal: \(-(-1)^k\)\\[4pt] Above diagonal: \((-1)^k + k(-1)^k\)\\[4pt] Top right: \(-k(-1)^k - \frac{k(k-1)}{2}(-1)^k\)} & M1 A1 & {\msNote{M1}{Multiplies their assumed matrix by the given matrix, finding all entries or the three distinct non-zero types. Condone sign slips.}\\ \msNote{A1}{All entries correct, need not be simplified.}} \\
& {\((-1)^{k+1}\), \(\;-(k+1)(-1)^{k+1}\)\\[4pt] \(\frac{(k+1)k}{2}(-1)^{k+1}\), zeros below the diagonal} & A1 & {cso. All entries shown in the given form with \(n=k+1\); top right as \(\frac{(k+1)k}{2}(-1)^{k+1}\) oe.} \\
& {True for \(n=1\); true for \(n=k\) \(\Rightarrow\) true for \(n=k+1\)\\[4pt] \(\Rightarrow\) true for all \(n \in \mathbb{Z}^+\)} & E1 & {Concludes with all three points: true for \(n=1\), true for \(n=k\) implies true for \(n=k+1\), so true for all positive integers \(n\).} \\
\nopagebreak \SetCell[c=2]{r} Total for Question 6 & & 6 & \\
\end{msscheme}

\msQuestion{7}
\begin{msscheme}{
  hline{3-6} = {2-Z}{0.4pt, msGrid},
  hline{7} = {0.4pt, msGrid},
  row{7} = {font=\bfseries},
}
Question & Scheme & Marks & Notes \\
7 & {\(F_2 = F_1 + F_0 = 1\)\\[4pt] \(\Rightarrow \mathbf{M}^1 = \begin{pmatrix}1&1\\1&0\end{pmatrix} = \begin{pmatrix}F_2&F_1\\F_1&F_0\end{pmatrix}\)} & B1 & {Needs \(F_2 = 1\) from the recurrence and \(\mathbf{M}\) matched entry by entry with \(F_2, F_1, F_0\).} \\
& {Assume \(\mathbf{M}^k = \begin{pmatrix}F_{k+1}&F_k\\F_k&F_{k-1}\end{pmatrix}\)\\[4pt] \(\mathbf{M}^{k+1} = \mathbf{M}^k\mathbf{M}\)\\[4pt] \(\phantom{\mathbf{M}^{k+1}} = \begin{pmatrix}F_{k+1}&F_k\\F_k&F_{k-1}\end{pmatrix}\begin{pmatrix}1&1\\1&0\end{pmatrix}\)} & M1 & {States the result for \(n=k\) and writes \(\mathbf{M}^{k+1}\) as \(\mathbf{M}^k\mathbf{M}\) or \(\mathbf{M}\mathbf{M}^k\) with the assumed form substituted.} \\
& {\(\phantom{\mathbf{M}^{k+1}} = \begin{pmatrix}F_{k+1}+F_k & F_{k+1}\\F_k+F_{k-1} & F_k\end{pmatrix}\)} & M1 A1 & {\msNote{M1}{Multiplies the two matrices row by column to four entries in terms of \(F\); condone one slip.}\\ \msNote{A1}{Correct product, simplified or not, from either order of multiplication.}} \\
& {\(F_{k+1}+F_k = F_{k+2}\), \(F_k+F_{k-1} = F_{k+1}\)\\[4pt] \(\Rightarrow \mathbf{M}^{k+1} = \begin{pmatrix}F_{k+2}&F_{k+1}\\F_{k+1}&F_k\end{pmatrix}\)} & A1 & {cso. Uses the recurrence on both sums to reach the \(n=k+1\) form.} \\
& {True for \(n=1\); true for \(n=k\)\\[4pt] \(\Rightarrow\) true for \(n=k+1\);\\[4pt] so true for all integers \(n \geq 1\)} & E1 & {States true for \(n=1\), that true for \(n=k\) implies true for \(n=k+1\), and hence true for all integers \(n \geq 1\).} \\
\nopagebreak \SetCell[c=2]{r} Total for Question 7 & & 6 & \\
\end{msscheme}

\msQuestion{8}
\begin{msscheme}{
  hline{3-6} = {2-Z}{0.4pt, msGrid},
  hline{7} = {0.4pt, msGrid},
  row{7} = {font=\bfseries},
}
Question & Scheme & Marks & Notes \\
8 & {\(n=1\): \(2^1-1=1\), \(1+(1-2)2^0=0\),\\[4pt] \(1\times 2^0=1\), \(2^1=2\) \(\Rightarrow\) formula gives \(\mathbf{A}\)} & B1 & {Formula evaluated at \(n=1\) and shown equal to \(\mathbf{A}\). Not for "true for \(n=1\)" alone.} \\
& {Assumes the formula true for \(n=k\)\\[4pt] \(\mathbf{A}^{k+1}=\mathbf{A}^k\mathbf{A}\), with the assumed \(\mathbf{A}^k\)} & M1 & {States the assumption for \(n=k\) and writes \(\mathbf{A}^{k+1}\) as \(\mathbf{A}^k\mathbf{A}\) or \(\mathbf{A}\mathbf{A}^k\) using it.} \\
& {Row 1: \(1,\ 1+2(2^k-1),\)\\[4pt] \(\phantom{Row 1: 1,\ }(2^k-1)+2\bigl(1+(k-2)2^{k-1}\bigr)\)\\[4pt] Row 2: \(0,\ 2^{k+1},\ 2^k+2k\,2^{k-1}\)\\[4pt] Row 3: \(0,\ 0,\ 2^{k+1}\)} & M1 A1 & {\msNote{M1}{Multiplies their \(\mathbf{A}^k\) and \(\mathbf{A}\) row by column, in either order, attempting every non-zero entry. Condone slips.}\\ \msNote{A1}{Correct product, unsimplified, oe for \(\mathbf{A}\mathbf{A}^k\).}} \\
& {Entries \(2^{k+1}-1\), \(1+(k-1)2^k\), \((k+1)2^k\), \(2^{k+1}\)\\[4pt] \(\Rightarrow\) the formula with \(n=k+1\)} & A1 & {cso. Every entry shown to match the formula with \(n=k+1\); \(1+(k-1)2^k\) and \((k+1)2^k\) need not be rewritten in terms of \(k+1\).} \\
& {True for \(n=1\); if true for \(n=k\), true for \(n=k+1\)\\[4pt] \(\Rightarrow\) true for all positive integers \(n\)} & E1 & {Concludes with all three: true for \(n=1\); true for \(n=k\) implies true for \(n=k+1\); hence true for all positive integers \(n\).} \\
\nopagebreak \SetCell[c=2]{r} Total for Question 8 & & 6 & \\
\end{msscheme}

\end{document}
